\section{Verifier 如何进入完整自我改进闭环}

\subsection{推理阶段}

\begin{itemize}
    \item Best-of-$N$：完整生成后 rerank；
    \item Beam / tree search：对前缀评分并剪枝；
    \item Early stopping：置信度足够时停止生成；
    \item Critique and revise：把低分原因反馈给 generator 修订。
\end{itemize}

\subsection{训练阶段}

\begin{itemize}
    \item 数据过滤：只保留高分 synthetic trajectory；
    \item Rejection sampling fine-tuning：用 verifier 选择的轨迹做 SFT；
    \item Reinforcement learning：把 verifier 分数作为 reward；
    \item Curriculum：按步骤难度和错误类型组织训练样本。
\end{itemize}

\subsection{闭环稳定性检查}

\begin{enumerate}
    \item Generator 和 verifier 是否使用独立评估集？
    \item Verifier 更新后，旧 reward 是否仍可比较？
    \item RL 是否提高 verifier 分数却降低真实正确率？
    \item 新 generator 是否产生 verifier 未见过的攻击性分布？
    \item 是否保留人工审计与外部 ground truth 作为锚点？
\end{enumerate}

\begin{importantbox}{Generator--Verifier 共进化必须有外部锚点}
如果 generator 只优化 verifier，而 verifier 又只从 generator 数据学习，两者可能共同漂移到一个自洽但错误的协议。可执行测试、形式化规则、人工抽检和保留 benchmark 是防止闭环失真的关键。
\end{importantbox}

\subsection{本章小结}

Verifier 是选择器、搜索启发式、数据过滤器和 reward model。它的偏差会沿闭环被放大，因此自我改进系统的首要工程任务不是“多跑几轮”，而是证明反馈信号仍与真实目标一致。

